\chapter{Conclusions}\label{ch:conclusion}
The numerical behavior of $(x^{10}+c)-x^{10}$ can be understood through the ULP of the reused stored power $y$ and the low-order binary structure of the stored increment $c$. Away from binade crossings and with $y\ge c$, the output equals the increment rounded to the nearest multiple of $\ulp(y)$, with ties determined by the parity of the significand of $y$. This produces an initial exact regime, a parity-sensitive first inexact binade, discrete rounded plateaus, and complete absorption.

Writing $c=m2^{\nu}$ with odd $m$ and $\lambda=\lfloor\log_2c\rfloor$, the model predicts $x_1\approx2^{(p+\nu)/10}$ and $x_2\approx2^{(p+1+\lambda)/10}$, with $L=\lambda-\nu+1$ potentially inexact binades. The distinction between one-bit increments and long-significand fractions explains why powers of two produce narrow transition bands whereas stored approximations to values such as $0.1$ or $1/3$ show error over a much broader range. In binary64, odd integers first become inexact near $x=39.40$; in binary32 the corresponding point is near $x=5.28$.

The supplied study reports a 1,611-increment validation in each precision, while the independently repeated 780-case power-of-two classification agrees with the preliminary report. Explicit assumptions and the exceptions at binade boundaries remain necessary. The example is pedagogically valuable because it exposes a key floating-point principle: an operation can be perfectly well rounded relative to its \emph{large} operand and still completely destroy the information in its \emph{small} operand.
